\section{Cierre y ampliación: del MoE de lenguaje al Vision MoE}

Al final de la clase, primero se resumen los resultados de Switch en modelado de lenguaje y después se presenta un trabajo nuevo en el ámbito visual. Las dos partes muestran que la interfaz central de los expertos dispersos puede reutilizarse entre modalidades, pero que el significado de un token, la estructura de redundancia y el orden de importancia cambian la estrategia de routing. En texto, por lo general se busca que cada token reciba un cómputo similar; con los patches de una imagen, en cambio, puede admitirse un procesamiento selectivo más agresivo.

\lecturefigure{slide-36-wrapping-up.jpg}{Resumen de la clase: enrutamiento simple, entrenamiento estable, aceleración del preentrenamiento, multilingüismo y escala de billones de parámetros.}{Video oficial de Stanford Online 00:57:08--01:00:43.}

\begin{knowledgebox}{Lectura de la figura: convertir los «resultados» en enunciados condicionales}
Switch simplifica el MoE y, dados una línea base de T5, los datos de C4/mT5, TPU y una implementación en Mesh TensorFlow, obtiene ganancias de velocidad y de calidad; entrena un modelo de un billón de parámetros, pero cada token activa solo unos pocos experts; mejora el desempeño multilingüe promedio, pero aún hay que revisar el enrutamiento por idioma (per-language) y el costo de despliegue.
\end{knowledgebox}

El MoE visual trata los image patches como tokens e inserta sparse expert layers en el Transformer. Las imágenes tienen mucha redundancia local, por lo que el router no solo puede decidir «qué expert usar», sino también «qué patches merecen ocupar la capacity limitada». Con esto, la sparsity empieza a acercarse a la adaptive computation.

\lecturefigure{slide-37-sparse-models-computer-vision.jpg}{V-MoE: extensión del sparse expert scaling a la clasificación de imágenes.}{Video oficial de Stanford Online 01:00:44--01:01:15.}

\begin{knowledgebox}{Qué es un token visual}
El Vision Transformer divide la imagen en patches y proyecta cada patch como un token. En V-MoE, lo que procesa cada expert es la representación de un patch; no se elige un solo expert para cada imagen completa. Patches distintos pueden enrutarse a experts distintos.
\end{knowledgebox}

Cuando el capacity factor es menor que 1, la priority routing primero estima la importancia de cada token y después asigna los espacios de los experts a los patches de mayor prioridad. Si $\gamma=0.5$, la capacidad alcanza solo para la mitad de los tokens de una distribución idealmente uniforme; ya no se trata solo de un overflow accidental, sino de una asignación deliberada del presupuesto de cómputo. El fondo de la imagen o las regiones repetidas pueden omitirse, y los patches de los objetos relevantes reciben más cómputo de las capas dispersas.

\lecturefigure{slide-38-priority-routing-vision-moe.jpg}{Priority routing en Vision MoE: cuando la capacidad no alcanza, se procesan primero los patches importantes.}{Video oficial de Stanford Online 01:01:16--01:02:10.}

\teachervoice{El ponente considera muy prometedor bajar el capacity factor por debajo de 1: en imágenes, incluso si solo cerca de una décima parte de los patches pasa por la sparse layer, la calidad puede mantenerse buena. En ese punto, la sparsity empieza a modificar a la vez la «selección de parámetros» y la «asignación del cómputo».}

\begin{warningbox}{El priority score también es una hipótesis aprendida}
Un patch omitido no necesariamente es «poco importante». Si el priority router asigna de forma sistemática puntajes bajos a objetos pequeños, texto, regiones de borde o clases minoritarias, la exactitud promedio puede ocultar esa fragilidad. Conviene hacer pruebas de oclusión, por subgrupos, de calibración y fuera de distribución.
\end{warningbox}

\subsection{Límites de despliegue señalados en las preguntas y respuestas}

La sesión de Q\&A del final añadió tres límites que las slides no presentan completos. Primero, con C4, SuperGLUE y las longitudes de secuencia habituales de entonces, el attention map no era el principal costo de memoria; los pesos pesaban más, así que la attention lineal no necesariamente resuelve primero el cuello de botella. Segundo, una de las restricciones directas para seguir aumentando los experts es que los parámetros tienen que almacenarse en algún lugar. Tercero, la eficiencia por parámetro almacenado de un sparse model suele ser menor que la de un dense model; por eso resulta más adecuado para el preentrenamiento y para servicios de rendimiento medio a alto, y no necesariamente para solicitudes individuales de baja latencia en los dispositivos más pequeños.

\begin{table}[H]
\centering
\small
\caption{Juicios por escenario surgidos en las preguntas y respuestas de la clase.}
\begin{tabularx}{\textwidth}{L{0.28\textwidth}X X}
\toprule
Escenario & Beneficio posible & Riesgo principal \\
\midrule
Preentrenamiento a gran escala & Los parámetros condicionales amplían la capacidad; alto rendimiento con paralelismo data/expert & all-to-all, almacenamiento, checkpoint y estabilidad \\
Servicio de rendimiento medio a alto & El lote puede agrupar tokens por expert y repartir el costo de comunicación & Experts saturados, deriva del enrutamiento y latencia entre máquinas \\
Solicitud individual de baja latencia / dispositivos de borde & Posible reducción de los active FLOPs & El total de pesos difícilmente cabe; un lote pequeño no permite formar GEMM de experts eficientes \\
Contexto muy largo & El MoE amplía la capacidad de la FFN & El costo de attention/KV puede volverse otro cuello de botella independiente \\
\bottomrule
\end{tabularx}
\end{table}

\teachervoice{El juicio final de Barret Zoph es mesurado: si el objetivo es «colocar el modelo más potente en el dispositivo más pequeño posible», la per-weight efficiency de un dense model puede ser mejor; si se trata de preentrenamiento masivo y de servicios de rendimiento medio a alto, es más probable que la sparsity aporte grandes beneficios. El escenario de aplicación determina la conclusión.}

\subsection{Resumen de la sección}

Los resultados de Switch en lenguaje muestran que los parámetros condicionales pueden escalarse; V-MoE va más allá y usa priority routing para repartir un cómputo limitado. Las preguntas y respuestas recuerdan que el almacenamiento, el lote, la interconexión, la longitud del contexto y la eficiencia por peso determinan si el sparse deployment ofrece de verdad una ventaja.
